\section{Discussion}
\label{sec:discussion}

\subsection{Crop-specific responses argue against a general spaceflight ionome}

The clearest result is negative in a useful way: there is no shared
flight-induced ionome response across the two Veggie crops with archived
elemental data. Lettuce shows no significant change in any element on either
scale; Tokyo Bekana shows \NSigCLR{} after compositional correction, and
\NDivergentCLR{} elements move in opposite directions between the crops. A
classifier trained on one crop performs below chance on the other.

This matters for how spaceflight crop nutrition is generalised. Statements of
the form ``spaceflight increases mineral $X$ in leafy greens'' are not supported
by these data even across two leafy green species grown in the same hardware.
The two studies differ in more than species --- mission, duration, light
spectrum, harvest protocol --- and those covariates cannot be separated with
$n=\NSamples$. That is precisely the point: with the available data, crop
identity and its confounds cannot be disentangled, so a species-general claim is
not currently derivable from the archive.

\subsection{Compositional artefacts are the dominant failure mode}

The apparent calcium and iron increases in Tokyo Bekana --- \FcTbCa$\times$ and
\FcTbFe$\times$, both large and one significant --- do not survive correction
for dry-matter dilution. Had we reported the raw concentrations, we would have
reported an accumulation effect that the data do not support.

This is a general hazard for ICP-OES tissue data, which is compositional by
construction \citep{aitchison1982}. Any treatment that changes tissue water
content, structural carbohydrate or non-mineral dry matter shifts every
concentration together. Spaceflight plausibly does all three. We would suggest
that reporting a concentration fold change without a compositional or
content-based check should be treated as an incomplete analysis in this
literature, not a stylistic preference.

The elements that survive are informative in their own right. Potassium, sodium
and sulfur are the mobile ions most closely tied to water relations and
osmotic adjustment, which is a coherent picture for plants in altered
root-zone water distribution --- and consistent with the reduced water delivery
recorded for flight pillows.

\subsection{Data quality is a first-order scientific problem}

The OSD-655 unit defect is a 1000-fold error on five of twelve elements, present
in both the data files and the archived data dictionary, in an openly available
and citable NASA dataset. It is not hidden --- a plausibility check catches it
immediately --- but nothing in the archive's own metadata flags it, and a
pipeline that trusts the stated unit inherits it silently.

We have written this up for the OSDR curation team
(\path{fair/osdr_curation_report.md}). The broader lesson is that a FAIR
dataset is interoperable only to the extent its unit metadata is correct, and
that automated plausibility gates belong in analysis pipelines as a matter of
routine. The gate that caught this is four lines of arithmetic.

\subsection{Overload is not the risk; sufficiency is the question}

The concern that motivated this work --- that plant-rich diets combined with
spaceflight-altered mineral handling could drive crew mineral overload --- is
not supported. Reaching NASA's iron upper limit would require
\OverloadFeLet\,kg of fresh lettuce daily, against a per-pillow yield of
roughly \BiomassFlight\,g. The margin is two orders of magnitude.

The useful question is the opposite one. A 100\,g serving supplies
\ServeLetK\% of daily potassium and \ServeLetCa\% of daily calcium: real but
supplementary contributions. The exception is manganese, at \ServeLetMn\% of
the daily requirement from a single serving, which is worth monitoring as
Veggie-type production scales --- not because these values approach a toxic
threshold, but because it is the one element where produce could plausibly
become a dominant dietary source.

\subsection{The missing measurements}

Two absences constrain everything above, and both are cheap to fix.

\emph{Dry-matter fraction.} OSD-655 archives fresh harvest mass and per-dry-gram
concentrations but no dry-matter fraction, so total mineral content per plant
cannot be computed and uptake efficiency is unavailable. Dry-matter fraction is
routinely measured during ICP-OES sample preparation --- the tissue must be
dried to be analysed. Archiving that number would convert a bounded estimate
into a measurement, and would also allow the dilution correction of
Section~\ref{sec:clr} to be tested directly rather than inferred.

\emph{Substrate elemental composition.} No Veggie study archives the elemental
composition of the arcillite substrate and its controlled-release fertiliser
charge. Without it, bioavailability --- the fraction of supplied mineral that
reaches edible tissue --- cannot be computed at all. Our water-normalised index
is a proxy for the mass-flow driver, not a bioavailability fraction, and we
name it accordingly.

\subsection{What the crew-side data can and cannot support}

The Inspiration4 panel gives real, openly available serum mineral measurements,
and the post-flight potassium decline in \AstrodownK{} of \AstroN{} crew is a
genuine observation. But four subjects on a \AstroDuration-day mission with no
in-flight sampling cannot establish a spaceflight effect, and the plant and crew
data describe different missions and different people.

The dataset that would close this loop exists: the Nutritional Status
Assessment (SMO-016E), with 27 astronauts across long-duration ISS missions,
in-flight sampling, and paired dietary intake and serum chemistry
\citep{zheng2023}. Our LSDA survey confirms it is request-gated rather than
open. Section~\ref{sec:availability} and the request package in
\path{lsda_request/} set out the route.

\subsection{Limitations}

\begin{itemize}\itemsep2pt
  \item $n=\NSamples$ plant samples in total (\NLet{} lettuce, \NTb{} Tokyo
        Bekana). The crop divergence needs replication before it is treated as
        established.
  \item Crop, mission, duration and light spectrum are confounded; the
        divergence is between studies, not cleanly between species.
  \item The compositional correction assumes the dilution factor is
        multiplicative and shared across elements within a sample. Direct
        dry-matter measurement would test this.
  \item Water and mass supplements cover VEG-04A only
        (\SupplyMatched/\SupplyTotal{} samples), and biomass is fresh weight.
  \item Serving calculations assume a $5$--$10\%$ dry-matter fraction from the
        literature, not from these studies.
  \item The astronaut cohort is $n=\AstroN$ over \AstroDuration{} days with no
        in-flight sample.
  \item LSDA annotation is text-matching over catalog metadata; an experiment
        whose record does not name its analytes will be missed, so
        \LSDAMeasured{} is a lower bound.
\end{itemize}
